\subsection{\texttt{TROWEXPAND}}
\noindent\textbf{Bundle encoding.} UOP entry 44, opcode \texttt{0x104D}. Bundle: \texttt{B.OP + B.ARG + B.DIM + B.IOTI}. \texttt{B.OP.Opcode}=\texttt{0x104D}. \texttt{B.IOTI}: selector=\texttt{111 last}, \texttt{SrcTile0}=\texttt{src}, \texttt{SrcTile1}=\texttt{absent}, \texttt{DstTile}=\texttt{dst}, \texttt{SizeCode}=profile tile size. \texttt{B.DIM} supplies row-axis reduction or row-broadcast bounds.\par\smallskip
{\small
\paragraph{PTO ISA description.}
Broadcast the first element of each source row across the destination row.
\paragraph{Example.}
Source column vector [[a], [b], [c]] (3x1) broadcast to 3x4: result row 0 = [a, a, a, a], row 1 = [b, b, b, b], row 2 = [c, c, c, c].
\par\smallskip
\paragraph{PTO ISA operation.}
Let \texttt{R = dst.GetValidRow()} and \texttt{C = dst.GetValidCol()}. For \texttt{0 \textless{}= i \textless{} R} and \texttt{0 \textless{}= j \textless{} C}:

\[
\mathrm{dst}_{i,j} = \mathrm{src}_{i,0}
\]
}\par\smallskip
\paragraph{Notes.}
Only the first column of the source tile is used. The j-th output row always uses src[j,0].
\paragraph{Microcode site table.}
{\scriptsize
\begin{longtable}{@{}R{0.055\linewidth}L{0.23\linewidth}L{0.31\linewidth}L{0.22\linewidth}@{}}
\toprule
Seq & Micro instruction & WO[63:32] fields & WM[31:0] fields \\
\midrule
0 & \texttt{u\_vec.vlds} & \texttt{opc=0x17, x=0, fl=mem, seq=0, kind=b} & \texttt{entry=44, cnt=2, res=vec, var=0} \\
1 & \texttt{u\_vec.vdup} & \texttt{opc=0x13, x=0, fl=alu, seq=1, kind=u} & \texttt{entry=44, cnt=1, res=vec, var=0} \\
2 & \texttt{u\_vec.vsts} & \texttt{opc=0x2A, x=0, fl=mem, seq=2, kind=b} & \texttt{entry=44, cnt=2, res=vec, var=0} \\
\bottomrule
\end{longtable}
}
\paragraph{Microcode body DSL.}
\begin{lstlisting}[style=ptocode]
def uop_trowexpand(src, dst):
    dtype = dst.element_type
    valid_rows, valid_cols = dst.valid_shape

    for row in range(0, valid_rows, 1):
        remained = valid_cols
        for col in range(0, valid_cols, pto.get_lanes(dtype)):
            mask, remained = pto.make_mask(dtype, remained)
            scalar_vec = pto.vlds(src[row, :])   # load entire row from src (broadcast semantics: first element replicated across lanes by vdup)
            broadcasted = pto.vdup(scalar_vec, mask)
            pto.vsts(broadcasted, dst[row, col:], mask)
    return
\end{lstlisting}
\Needspace{0.34\textheight}
